% Part: computability
% Chapter: recursive-functions
% Section: other-recursions

\documentclass[../../../include/open-logic-section]{subfiles}

\begin{document}

\olfileid{cmp}{rec}{ore}
\olsection{पुनरावर्तनाचे इतर प्रकार}

जोडीकरण आणि क्रमिका वापरून आदिम पुनरावर्तनाच्या अधिक असामान्य (आणि
उपयोगी) रूपांचे समर्थन करता येते. उदाहरणार्थ, पुढील व्याख्येप्रमाणे दोन
फलने एकसामयिकपणे परिभाषित करणे अनेकदा उपयोगी ठरते:
\begin{align*}
h_0(\vec x, 0) & = f_0(\vec x) \\
h_1(\vec x, 0) & = f_1(\vec x) \\
h_0(\vec x, y+1) & = g_0(\vec x, y, h_0(\vec x, y), h_1(\vec x, y)) \\
h_1(\vec x, y+1) & = g_1(\vec x, y, h_0(\vec x, y), h_1(\vec x, y))
\end{align*}
हे \emph{एकसामयिक पुनरावर्तनाचे} एक उदाहरण आहे. फलने परिभाषित करण्याची
आणखी एक उपयोगी पद्धत म्हणजे $h(\vec x, y+1)$ चे मूल्य आधीच्या
$h(\vec x, 0)$, \dots,~$h(\vec x, y)$ या \emph{सर्व} मूल्यांच्या आधारे
देणे; उदाहरणार्थ:
\begin{align*}
h(\vec x, 0) & = f(\vec x) \\
h(\vec x, y+1) & = g(\vec x, y, \tuple{h(\vec x, 0), \dots, h(\vec x, y)}).
\end{align*}
पुढील योजना ही कल्पना अधिक संक्षिप्तपणे व्यक्त करते:
\[
h(\vec x, y) = g(\vec x, y, \tuple{h(\vec x, 0), \dots, h(\vec x, y-1)})
\]
येथे $y$ हे $0$ असेल, तेव्हा $g$ चा शेवटचा आदान केवळ रिकामी क्रमिका
असतो, असे समजायचे. दोन्ही मांडण्यांमागील कल्पना अशी की, ``उत्तरवर्ती
टप्प्याचे'' संगणन करताना $h$ ला आतापर्यंत संगणित केलेल्या मूल्यांची संपूर्ण
क्रमिका वापरता येते. याला \emph{मूल्यक्रम पुनरावर्तन} म्हणतात. एका विशिष्ट
उदाहरणात, पुढील प्रकारच्या व्याख्येचे समर्थन करण्यासाठी ते वापरता येते:
\begin{align*}
h(\vec x, y) & = \begin{cases}
  g(\vec x, y, h(\vec x, k(\vec x, y))) & \text{जर $k(\vec x, y) < y$ असेल} \\
  f(\vec x) & \text{अन्यथा.}
\end{cases}
\end{align*}
दुसऱ्या शब्दांत, $y$ येथील $h$ चे मूल्य~$k$ ने दिलेल्या \emph{कोणत्याही}
आधीच्या मूल्यावरील $h$ च्या मूल्याच्या आधारे संगणित करता येते.


\begin{prob}
  शेष फल $r(x,y)$ ची मूल्यक्रम पुनरावर्तनाने व्याख्या करा. ($x$, $y$ या
  नैसर्गिक संख्या असतील आणि $y > 0$ असेल, तर $r(x,y)$ ही~$y$ पेक्षा लहान
  अशी संख्या असते की, एखाद्या~$z$ साठी $x = z\times y + r(x,y)$. निश्चिततेसाठी,
  $y=0$ असल्यास $r(x,0) = 0$ असे मानू.)
\end{prob}

ही फलने नेहमीच्या आदिम पुनरावर्तनाने कशी मिळवायची याचा विचार करा. आदिम
पुनरावर्तनाची आणखी एक शेवटची आवृत्ती अधिक लवचीक आहे: तिच्यात वाटेत
\emph{प्राचल} (बाजूची मूल्ये) बदलण्याची मुभा असते:
\begin{align*}
h(\vec x, 0) & = f(\vec x) \\
h(\vec x, y+1) & = g(\vec x, y, h(k(\vec x), y))
\end{align*}
याचाही नेहमीच्या आदिम पुनरावर्तनाने अनुकार करता येतो. (हे करणे अवघड आहे.
सूचनेसाठी, संगणन हाताने उलट क्रमाने उलगडून पाहा.)

\end{document}
